\documentclass{beamer}
%\usepackage{fullpage}
%\usepackage[left=2.8cm,right=2.2cm,top=2 cm,bottom=2 cm]{geometry}
\setbeamersize{text margin left=10pt,text margin right=10pt}
\usepackage{amsmath,amssymb} 
\usepackage[T1]{fontenc}
\usepackage[utf8]{inputenc}
\usepackage[english,french]{babel}
\usepackage{txfonts}
\usepackage[]{graphicx}
\usepackage{multirow}
\usepackage{hyperref}

\hypersetup{
	colorlinks,
	urlcolor = blue
}

%\renewcommand{\baselinestretch}{1.5}

\def\supit#1{\raisebox{0.8ex}{\small\it #1}\hspace{0.05em}}

\title[BWEB: 01- Environnement] %
{Bureautique et Web \\Chapitre 01: Environnement de travail} 
\institute{ %
École  nationale Supérieure d'Informatique (ESI, ex. INI), Algérie
}
\author[ \textbf{\footnotesize  \insertframenumber/\inserttotalframenumber} \hspace*{\fill} ESI (2019-2020)] %
{ARIES Abdelkrime}
%\titlegraphic{\includegraphics[height=1cm]{../img/esi-logo.png}%\hspace*{4.75cm}~


\date{Année unniversitaire: 2019/2020} %\today

\usetheme{Warsaw} % Antibes Boadilla Warsaw

\beamertemplatenavigationsymbolsempty

%\setbeamertemplate{headline}{}


\begin{document}

\begin{frame}[plain]
	\maketitle
\end{frame}

\section{Introduction aux environnements de travail}

\begin{frame}
\frametitle{Introduction aux environnements de travail}

Un environnement de travail doit avoir:
\begin{itemize}
	\item un poste de travail: l'ordinateur et ses périphériques
	\item des applications: 
	\item espaces de stockage
\end{itemize}

\end{frame}

\subsection{Composition d'un environnement de travail}

\begin{frame}
\frametitle{Introduction aux environnements de travail}
\framesubtitle{Composition d'un environnement de travail}
%TODO schéma
\begin{itemize}
	\item Processeur 
	\item RAM 
	\item Stockage
	\item Périphériques 
	\item Système d'exploitation
	\item Applications 
\end{itemize}

\end{frame}

\subsection{Types d'environnements de travail}

\begin{frame}
\frametitle{Introduction aux environnements de travail}
\framesubtitle{Types d'environnements de travail}

Selon l'emplacement des données et des applications, un environnement de travail peut être: 
\begin{itemize}
	\item local
	\item distant
\end{itemize}

\end{frame}

\section{Environnement de travail local}

\begin{frame}
\frametitle{Environnement de travail local}

\end{frame}

\subsection{Matériel informatique}

\begin{frame}
\frametitle{Environnement de travail local}
\framesubtitle{Matériel informatique}

\end{frame}

\subsection{Système d'exploitation}

\begin{frame}
\frametitle{Environnement de travail local}
\framesubtitle{Système d'exploitation}

Parmi les fonctions d'un système d'exploitation: 
\begin{itemize}
	\item 
\end{itemize}

\end{frame}

\begin{frame}
\frametitle{Environnement de travail local}
\framesubtitle{Système d'exploitation}

SE de bureau (Desktop OS):
\begin{itemize}
	\item Windows (Microsoft)
	\item MacOS (Apple)
	\item Linux (Linux foundation)
	\item Chrome OS (Google): basé sur Linux
	\item BSD
\end{itemize}

SE de mobiles (Mobile OS): 
\begin{itemize}
	\item Android (Google)
	\item iOS (Apple)
	\item Windows 10 Mobile (Microsoft) 
	\item Chrome OS (Google)
	\item BSD
\end{itemize}

\end{frame}

\begin{frame}
\frametitle{Environnement de travail local}
\framesubtitle{Système d'exploitation}

%D'après \url{https://gs.statcounter.com/os-market-share/desktop/worldwide/#monthly-201901-201911}
% https://www.w3schools.com/browsers/browsers_os.asp

%\includegraphics[width=\textwidth]{imagefile}

\end{frame}

\begin{frame}
\frametitle{Environnement de travail local: Windows}
\framesubtitle{Système d'exploitation: Windows}

\end{frame}

\begin{frame}
\frametitle{Environnement de travail local: Windows}
\framesubtitle{Système d'exploitation: Windows}

\end{frame}

\begin{frame}
\frametitle{Environnement de travail local: Linux}

\end{frame}

\begin{frame}
\frametitle{Environnement de travail local: Linux}
\framesubtitle{Distributions}

Les distributions majeures d'après DistroWatch: % https://distrowatch.com/dwres.php?resource=major

\begin{tabular}{llll}
	\includegraphics[width=.2\textwidth]{..//img/Bweb01-environnement/mint-logo.png} & \includegraphics[width=.2\textwidth]{..//img/Bweb01-environnement/ubuntu-logo.png} & 
	\includegraphics[width=.2\textwidth]{..//img/Bweb01-environnement/mx-logo.png} &
	\includegraphics[width=.2\textwidth]{..//img/Bweb01-environnement/arch-logo.png} \\
	
	\includegraphics[width=.2\textwidth]{..//img/Bweb01-environnement/gentoo-logo.png} & 
	\includegraphics[width=.2\textwidth]{..//img/Bweb01-environnement/slackware-logo.png} &
	\includegraphics[width=.2\textwidth]{..//img/Bweb01-environnement/debian-logo.png} & 
	\includegraphics[width=.2\textwidth]{..//img/Bweb01-environnement/fedora-logo.png} \\
	
\end{tabular}

\end{frame}

\begin{frame}
\frametitle{Environnement de travail local: Linux}
\framesubtitle{Environnements de bureau}

\end{frame}





%\subsection{Bibliography}
%\frame[allowframebreaks]%
%{\frametitle{Bibliography}
%\tiny
%\bibliography{biblio}
%\bibliographystyle{apalike} 
%}


\end{document}

